\section{Tutorial 11: CNN and Sequence Modeling Intuition}

\subsection{Objectives}

This tutorial develops intuition for structure-aware neural networks:

\begin{itemize}
    \item Understand how CNNs process spatial data
    \item Understand how RNNs process sequential data
    \item Compare architectural inductive biases
    \item Relate structure to efficiency and performance
\end{itemize}

\medskip

\noindent
\textbf{Focus.}  
How architectural design reflects the structure of data.

\subsection{Exercise 1: Fully Connected vs Convolution}

\textbf{Question.}  
Why is a fully connected network inefficient for image data?

\medskip

\noindent
\textbf{Discussion.}

\begin{itemize}
    \item ignores spatial locality
    \item requires many parameters
    \item does not exploit repeating patterns
\end{itemize}

\medskip

\noindent
\textbf{Insight.}  
Ignoring structure leads to inefficiency and poor generalization.

\subsection{Exercise 2: Local Receptive Fields}

Consider a convolutional filter applied to an image.

\medskip

\noindent
\textbf{Tasks.}
\begin{enumerate}
    \item What does a $3 \times 3$ filter capture?
    \item How does stacking layers increase receptive field?
\end{enumerate}

\medskip

\noindent
\textbf{Solution.}

\begin{itemize}
    \item $3 \times 3$ filter captures local patterns (edges, textures)
    \item stacking layers aggregates larger spatial context
\end{itemize}

\medskip

\noindent
\textbf{Insight.}  
CNNs build global understanding from local computations.

\subsection{Exercise 3: Weight Sharing}

\textbf{Question.}  
Why is weight sharing important in CNNs?

\medskip

\noindent
\textbf{Discussion.}

\begin{itemize}
    \item same feature can appear anywhere
    \item reduces number of parameters
    \item improves generalization
\end{itemize}

\medskip

\noindent
\textbf{Insight.}  
Weight sharing encodes translation symmetry.

\subsection{Exercise 4: Hierarchical Features}

\textbf{Task.}  
Describe what different CNN layers learn.

\medskip

\noindent
\textbf{Discussion.}

\begin{itemize}
    \item early layers: edges and textures
    \item middle layers: shapes and parts
    \item deep layers: objects and semantics
\end{itemize}

\medskip

\noindent
\textbf{Insight.}  
Deep networks build hierarchical representations.

\subsection{Exercise 5: Sequential Data Processing}

Consider a sequence $(x_1, x_2, \dots, x_T)$.

\medskip

\noindent
\textbf{Task.}  
Explain how an RNN processes this sequence.

\medskip

\noindent
\textbf{Solution.}

\begin{itemize}
    \item processes inputs one step at a time
    \item maintains hidden state $h_t$
    \item summarizes past information in $h_t$
\end{itemize}

\medskip

\noindent
\textbf{Insight.}  
RNNs compress past information into a state.

\subsection{Exercise 6: Information Bottleneck in RNNs}

\textbf{Question.}  
Why is it difficult for RNNs to capture long-term dependencies?

\medskip

\noindent
\textbf{Discussion.}

\begin{itemize}
    \item all past information must pass through $h_t$
    \item gradients degrade over long sequences
\end{itemize}

\medskip

\noindent
\textbf{Insight.}  
The hidden state acts as a bottleneck.

\subsection{Exercise 7: CNN vs RNN Comparison}

Fill in the table:

\[
\begin{array}{c|c|c}
\text{Aspect} & \text{CNN} & \text{RNN} \\
\hline
\text{Structure} & & \\
\text{Computation} & & \\
\text{Parallelism} & & \\
\text{Inductive bias} & & \\
\end{array}
\]

\medskip

\noindent
\textbf{Discussion.}

\begin{itemize}
    \item CNN:
    \begin{itemize}
        \item spatial structure
        \item parallel computation
        \item locality + translation invariance
    \end{itemize}

    \item RNN:
    \begin{itemize}
        \item temporal structure
        \item sequential computation
        \item memory over time
    \end{itemize}
\end{itemize}

\medskip

\noindent
\textbf{Insight.}  
Different architectures encode different structural assumptions.

\subsection{Exercise 8: When to Use Which?}

\textbf{Task.}  
Match model types to data:

\begin{itemize}
    \item images
    \item language
    \item time series
\end{itemize}

\medskip

\noindent
\textbf{Solution.}

\begin{itemize}
    \item images $\rightarrow$ CNN
    \item language $\rightarrow$ RNN / Transformer
    \item time series $\rightarrow$ RNN / hybrid models
\end{itemize}

\medskip

\noindent
\textbf{Insight.}  
Model choice should reflect data structure.

\subsection{Exercise 9 (Discussion): Limitations}

\textbf{Question.}  
What are the limitations of CNNs and RNNs?

\medskip

\noindent
\textbf{Discussion.}

\begin{itemize}
    \item CNN:
    \begin{itemize}
        \item limited long-range interaction
    \end{itemize}

    \item RNN:
    \begin{itemize}
        \item sequential bottleneck
        \item difficult to parallelize
    \end{itemize}
\end{itemize}

\medskip

\noindent
\textbf{Insight.}  
Limitations motivate new architectures.

\subsection{Summary}

\begin{itemize}
    \item CNNs exploit spatial locality and symmetry
    \item RNNs model temporal dependencies via hidden states
    \item Both rely on structural inductive bias
    \item Limitations lead to attention-based models
\end{itemize}

\medskip

\noindent
\textbf{Preparation for next lecture.}  
We next study attention mechanisms, which overcome limitations of both CNNs and RNNs by modeling global interactions directly.